\par\bigskip\Needspace{6\baselineskip}
\pdfbookmark[0]{Bemerkungen zur Analysis Residuorum}{bemerkungen-analysis-residuorum}
\begin{center}\large\bfseries BEMERKUNGEN ZUR ANALYSIS RESIDUORUM.\end{center}
\BookRule{.09\linewidth}
\renewcommand{\BookRunningTitle}{Historische Herausgeberbemerkungen zur Analysis Residuorum}

Die beiden vorstehenden Abhandlungen sind einem umfangreichen Manuskript mit dem Titel Analysis Residuorum entnommen, das vermutlich aus dem Jahr 1797 oder 1798 stammt; durch eine vollständige Umarbeitung entstanden daraus später die Disquisitiones Arithmeticae. Der vollständige Titel des sechsten Kapitels lautet:

Lösung der Kongruenz $x^m-1\equiv0$ und der Gleichung $x^m-1=0$; mit Erläuterungen zur Theorie der regelmäßigen Vielecke.

Der wesentliche Inhalt seines zweiten Teils (\textsection\textsection.\ 253--278) ging in den siebten Abschnitt der Disqq. Arithm. über.

Außerdem ist das siebte Kapitel teilweise erhalten: Verschiedene Anwendungen einiger vorangehender Untersuchungen (\textsection\textsection.\ 279--302). Es gliedert sich in folgende Unterabschnitte:

Über Umformungen gewöhnlicher Brüche (\textsection\textsection.\ 279--281).

Über die Umwandlung gewöhnlicher Brüche in Dezimalbrüche (\textsection\textsection.\ 282--292).

Über die Lösung der unbestimmten Gleichung $xx=a+by$ (\textsection\textsection.\ 293--297).

Über die Lösung der unbestimmten Gleichung $axx+byy=c$ (\textsection\textsection.\ 298--301).

Über die Bestimmung der Teiler von Zahlen (\textsection.\ 302; die folgenden Bogen fehlen).

Dies alles ist fast wörtlich in den sechsten Abschnitt der Disqq. Arithm. aufgenommen.

Die beiden hier mitgeteilten Abschnitte behandeln die Gegenstände, die, wie aus der Vorrede und den Artikeln 11, 44, 61, 62, 65, 84 der Disqq. Arithm. hervorgeht, den Inhalt des achten Abschnitts dieses Werks bilden sollten. Es ist jedoch zu bemerken, dass dieser Plan später wieder geändert wurde; unter den Manuskripten findet sich nämlich ein Fragment mit der Überschrift Achter Abschnitt: Ausführlichere Betrachtung einiger Untersuchungen zur Kreisteilung. Dieses
\SourceLeaf{251}
beginnt mit Art.\ 367 und sollte somit die Fortsetzung der Disqq.\ Arithm.\ bilden; der Inhalt der wenigen noch vorhandenen Artikel ging aber später in die Abhandlung Summatio quarundam serierum singularium über, und deshalb wird dieses Fragment von der gegenwärtigen Ausgabe ausgeschlossen.

Im vorstehenden Abdruck der beiden Teile der Analysis Residuorum ist der Text des Originals im Wesentlichen treu beibehalten, obwohl es seiner Form nach nicht als druckfertig gelten kann; die folgenden Bemerkungen geben die wichtigsten Abänderungen an und fügen zugleich einige Erläuterungen hinzu.

\textsection.\ 237.\ Vgl.\ Disqq.\ Arithm.\ artt.\ 61, 62.

\textsection.\ 239.\ Vgl.\ Disqq.\ Arithm.\ artt.\ 53, 54, 65.

\textsection.\ 241.\ Ist $n=2^\nu$ und $\nu\geqq3$, so gibt es zwar keine Zahl $\rho$ der angegebenen Art, aber die gesamte Untersuchung wird dadurch nicht wesentlich geändert.

\textsection.\ 251.\ Vermutlich sollte die hier erwähnte Schwierigkeit durch die Einführung höherer Potenzen von $p$ als Moduln beseitigt werden. Vgl.\ \textsection\textsection.\ 363, 372, 373.

\textsection.\ 332.\ Die Voraussetzung, dass der Modul eine Primzahl ist, wird bis einschließlich \textsection.\ 372 beibehalten.

\textsection.\ 338.\ Das unvollständige Zitat kann sich auf Disqq.\ Arithm.\ art.\ 44 beziehen.

\textsection\textsection.\ 344--346.\ Von den beiden im Manuskript vorhandenen Beweisen ist hier der erste, der mit den Worten »wir gehen nun daran zu beweisen« eingeleitet wird und auf einer näheren Untersuchung der Ausdrücke $(1^\alpha2^\beta3^\gamma\ldots)$ beruht, nach der eigenen Anweisung des Verfassers vollständig unterdrückt (»Der gesamte vorangehende Beweis muss zusammen mit einem anderen Beweis des vorigen Satzes, den wir hinzufügen wollten, unterdrückt werden, da wir einen anderen, unendlich einfacheren entdeckt haben. Er beruht auf dieser Grundlage«); im obigen Abdruck ist ferner der zweite Beweis dadurch verkürzt, dass die Entwicklung von $\frac{x\,\mathrm{d}P}{P\,\mathrm{d}x}$ statt derjenigen von $\frac{x\,\mathrm{d}P}{\mathrm{d}x}$ betrachtet wird, wodurch zugleich ein im Original enthaltener Bezug auf den unterdrückten ersten Beweis vermieden wird.

\textsection.\ 348.\ Der Ausdruck radix prima ist hier in derselben Bedeutung zu verstehen wie radix propria in der Abhandlung Summatio quarundam serierum singularum, art.\ 11. --- Bei der Behauptung, dass die Koeffizienten $A'.B'\ldots$ des ausmultiplizierten Produkts ganze rationale Zahlen sind, wird auf das sechste Kapitel verwiesen, in dem aber die Theorie der Gleichung $x^\tau-1=0$ nur für den Fall behandelt wird, dass $\tau$ eine Primzahl ist; die Form des Beweises in \textsection.\ 349 führt zunächst zu folgender Ergänzung. Bringt man das ausmultiplizierte Produkt in die (für alle Wurzeln der Gleichung $\theta^\tau=1$ geltende) Form
\[
S=E+F\theta+\ldots+N\theta^{\tau-1}
\]
so sind die Koeffizienten $E,F\ldots N$ Polynome in $x$ mit ganzzahligen rationalen Koeffizienten. Da ferner das Produkt unverändert bleibt, wenn $\theta$ durch $\theta^k$ ersetzt wird, wobei $k$ eine beliebige zu $\tau$ teilerfremde Zahl bedeutet, gilt dasselbe für den Ausdruck $S$. Daraus folgt ohne Schwierigkeit, dass alle in $S$ enthaltenen Potenzen von $\theta$, deren Exponenten $s$ denselben größten gemeinsamen Teiler mit $\tau$ haben, auch identische Koeffizienten haben müssen. Da schließlich jede Summe solcher Potenzen $\theta^s$ stets eine ganze Zahl ist, ist offensichtlich, dass der Ausdruck $S$ und folglich auch das fragliche Produkt ein Polynom in $x$ mit ganzzahligen Koeffizienten ist, was zu zeigen war. Aus dieser Betrachtung ergibt sich zugleich die Richtigkeit der Bemerkung am Schluss des Paragraphen. Andere Gründe lassen jedoch vermuten, dass dem Verfasser schon damals der allgemeine Satz über die Transformation symmetrischer Funktionen (Demonstratio nova altera theorematis omnem functionem etc.\ art.\ 4) bekannt war, aus dem die obigen Sätze unmittelbar folgen.

\textsection.\ 352.\ Das Zeichen $R\equiv S(\mathrm{mod.}\,P)$ oder auch $R\equiv S(\mathrm{mod.}\,P,p)$ bedeutet hier und im
\SourceLeaf{252}
Folgenden, dass die Differenz $R-S$ modulo $p$ den Teiler $P$ hat. --- Das unvollständige Zitat kann sich auf Disqq.\ Arithm.\ art.\ 49 beziehen.

\textsection.\ 354.\ Durch Multiplikation mit $x^\nu-1$ ergibt sich, dass die Summen gleich hoher Potenzen der Wurzeln der beiden Gleichungen $(P,\rho^{k\nu+t})=0$, $(P,\rho^t)=0$ einander kongruent sind $(\mathrm{mod.}\,p)$, und daraus folgt die Kongruenz $(P,\rho^{k\nu+t})\equiv(P,\rho^t)(\mathrm{mod.}\,p)$, sobald $m<p$ gilt (vgl.\ \textsection.\ 244). Ist aber $m\geqq p$, so lässt sich der Koeffizient der Potenz $x^{m-p}$ in einer Gleichung nicht mehr modulo $p$ aus den gegebenen Potenzsummen ihrer Wurzeln bestimmen, weil er in den dazu dienenden Formeln von \textsc{Newton} mit dem Faktor $p$ auftritt. Tatsächlich darf man aus der Kongruenz je zweier gleich hoher Potenzsummen der Wurzeln der Gleichungen $A=0$, $B=0$ allgemein nur folgern, dass $A\equiv\mathfrak A^p\mathfrak C$, $B=\mathfrak B^p\mathfrak C(\mathrm{mod.}\,p)$ gilt, wobei $\mathfrak C$ den größten gemeinsamen Teiler der beiden Funktionen $A,B$ nach dem Primzahlmodul $p$ bezeichnet, während $\mathfrak A$ und $\mathfrak B$ völlig unbestimmte Funktionen sind. Es ist zu vermuten, dass der Verfasser die Allgemeingültigkeit des Satzes aus der Theorie der Transformation symmetrischer Funktionen und insbesondere aus dem folgenden Satz abgeleitet hat: Ist in Bezug auf einen beliebigen Modul $p$ die Differenz $R(x)-S(x)$ durch die Funktion $P(x)$ teilbar, und sind $a,b,c\ldots$ die Wurzeln der Gleichung $P(x)=0$, so sind die Funktionen
\[
\begin{gathered}
(x-R(a))(x-R(b))(x-R(c))\ldots\quad\text{und}\\
(x-S(a))(x-S(b))(x-S(c))\ldots
\end{gathered}
\]
einander modulo $p$ kongruent.

\textsection.\ 355.\ In \textsection.\ 368 wird gezeigt, dass $P$ und $\frac{\mathrm{d}P}{\mathrm{d}x}$ keinen gemeinsamen Teiler haben, wenn $P$ keinen Faktor mehr als einmal enthält.

\textsection\textsection.\ 358, 359.\ Die unter den Text gesetzte Note ist einem einzelnen Blatt entnommen, das wahrscheinlich den bereits in der Handschrift gestrichenen \textsection.\ 359 ersetzen sollte.

\textsection.\ 360.\ In der Formulierung des Satzes ist eine Ungenauigkeit der Handschrift berichtigt.

\textsection.\ 361.\ Hier bedeutet der Exponent $\frac1k$ im Zeichen $(\xi,\rho^{\frac1k})$ jede positive ganze Zahl $k'$, für die $kk'\equiv1(\mathrm{mod.}\,\nu)$ gilt, wobei $\nu$ die kleinste positive ganze Zahl ist, für die $x^\nu-1$ durch $\xi$ modulo $p$ teilbar wird; dabei ist vorauszusetzen, dass $\xi$ nicht durch $x$ modulo $p$ teilbar ist, und außerdem, dass $k$ zu $\nu$ teilerfremd ist. Die Richtigkeit der Behauptung, dass $\xi'$ durch $(\xi,\rho^{\frac1k})$ teilbar ist $(\mathrm{mod.}\,p)$, ergibt sich aus \textsection.\ 354.

\textsection.\ 363.\ Die Schlussbemerkung bezieht sich vermutlich auf die Einführung von Moduln, die Potenzen der Primzahl $p$ sind; vgl.\ \textsection\textsection.\ 251, 372, 373.

\textsection.\ 367.\ Die Wurzeln der Gleichung $x^3+xx-2x-1=0$ sind die zweigliedrigen Perioden, in die die Wurzeln der Gleichung $\frac{x^7-1}{x-1}=0$ zerfallen. Dasselbe Beispiel findet sich auch auf einem einzelnen Blatt, auf dem das Hauptresultat der \textsection\textsection.\ 362, 363 unter dem Titel »der goldene Lehrsatz« ausgesprochen ist.

\textsection.\ 371.\ Dieser Paragraph sollte ein Beispiel enthalten; es ist jedoch nicht ausgeführt.
\begin{flushright}R.\ \textsc{Dedekind}.\end{flushright}
\BookRule{.28\linewidth}
